\section{Introduction}

The standard agent loop pays for a model decision at every tool
boundary~\citep{yao2023react}. That flexibility is valuable while a workflow is
novel, but mature agents often revisit the same read-only evidence path with the
same data dependencies, where repeating the decision adds latency, tokens, and
cost without adding adaptation.

Removing that decision is not ordinary caching: a tool argument may depend on
an earlier observation, a repeated sequence may cross an approval or write, a
read-only tool may hide an effect, and replaying the same outputs may still
change the answer. The question is when a recurrent model-mediated region may
be replaced by a deterministic program, and when it must be refused.

Consider held-out issue \#4420. Given only \code{issue\_number=4420}, the
observed trace calls \code{record(4420)}, passes the \emph{returned}
\code{record.issue\_number} to \code{labels} and \code{comments}, and invokes
the latter with \code{limit=100}. All 132 discovery executions take that
three-read route (130 with \code{limit=100}, two with \code{limit=1}), so a
recurrence-only optimizer would emit all three reads. \method instead
retires the \emph{region}: \code{comments.limit} is not an invariant literal
and is reachable neither from the entry state nor from either earlier result,
so one unwitnessed slot blocks the three-read candidate as
\code{ungroundable\_slot}. What survives is the maximal
justified prefix \code{record} $\rightarrow$ \code{labels}; the agent still
chooses the limit and writes the answer. That two-read artifact passes 92/92
calibration groups at $\alpha=0.05$ (upper bound 0.0498).

An earlier 18-record study supplies the counterexample: its three-read
artifact, admitted at a looser $\alpha=0.10$ (a 45-group floor), passed all 45
calibration replays and the tool contract on every
held-out record, yet issue \#6602 lost the URL from a Markdown link and only
17/18 downstream answers stayed exact; a provider-free continuation check, run
after the model's answer, detects that miss. Together the two cohorts show why
recurrence only proposes a candidate (\cref{fig:provenance} draws the first).


The closest precedent is trace compilation: Dynamo~\citep{bala2000dynamo}
and tracing JITs~\citep{gal2009tracejit} compile a hot path behind
\emph{guards} and fall back when one fails, as
deoptimization~\citep{holzle1992deopt} restores unoptimized execution.
Agent traces admit the same structure (family mining, an entry contract,
fallback to the unchanged agent) but compile only \emph{initial} read-only
prefixes (\cref{eq:prefix}) and add three obligations: tool arguments must be
shown \emph{grounded} in observable state; declared \emph{effects} must permit
speculation, since the ``interpreter'' calls the outside world; and the
residual error rate, being statistical, needs a finite-sample bound.

We study these obligations through \emph{guarded agentic compaction} (\method),
a compile-or-retire trace-to-program system for recurrent read-only prefixes
(\cref{fig:pipeline}). The core contribution is not ``agent compilation'' in
general, but an admissibility argument for trace-derived specialization:

\begin{enumerate}[leftmargin=1.2em,itemsep=1pt,topsep=2pt]
  \item A guarded specialization pipeline (\cref{alg:compile}) combining typed
provenance, effect and position barriers, bounded synthesis, runtime
verification, and finite-sample compile-or-retire admission
(\cref{alg:calibrate}, \cref{prop:admission}) with the certificate level and
its event stated per artifact. The bound is a standard exact binomial interval
under a union bound over a fixed grid; the contribution is the admission rule
and the composition that precedes it, to our knowledge the first such
certificate applied to trace-derived agent compilation. Every stage can refuse,
and the default output is the unchanged agent. \item A real-provider evaluation
across three public-record workflow families; a cross-repository extension on
newer records with admissions and principled retirements; a post-cutoff cohort
with an end-to-end certificate; a live SQL agent on the public BIRD benchmark
whose fixed schema-reading opening compiles on 26 of 31 databases
(\cref{sec:res-bird}); three
pre-registered drift ablations (two uninformative by construction, one adverse);
and replications on a second model family and provider.
  \item Calibrated refusal on public trace benchmarks (\cref{sec:res-refusal}):
  NESTFUL, the largest trace corpus the compiler ingests here, clears
  provenance, synthesis, and replay with no wrong execution yet retires below
  the exact sample floor (API-Bank and BFCL v4 likewise), a corpus-size fact
  that bounds what such benchmarks can certify; AppWorld admits a region that is
  eligible on 2{,}339/2{,}340 released full-code trajectories but 1/4{,}680 ReAct
  or plan-and-execute ones, separating \emph{admissible} from \emph{dispatchable}.
\end{enumerate}
